\begin{abstract}
We audit a prototype that was originally described as an agentic-fraud
detector.  Its real-data labels do not encode fraud, autonomous execution, or
human control.  The positive class contains 665 historical ERC-8004 registry
owner/default-wallet proxies and the comparison class contains 1,069 active
counterparties of unknown type.  On the retained local snapshot, the original
composite ranks these classes at ROC-AUC 0.5147.  A full-sample, post-hoc
reweighting reaches 0.6223; five-fold random-address cross-fitting gives pooled
out-of-fold AUC 0.6122, and a logistic combiner gives 0.6503.  These corrected
results are provisional until preserved in an immutable end-to-end run
manifest; only the historical 0.5147 is reproduced by the preserved audit.
They also remain exploratory because feature design, graph construction, cohort
cleaning, and model selection were informed by the same connected population.
At a training-fold threshold chosen for a nominal 5\% comparison-class alert
rate, held-out proxy recall is 7.5\% and proxy precision is 48.1\%.

The audit finds two decisive provenance defects.  First, a historical bridge
deduplicated ERC-20 events by transaction hash and discarded 13,539 of 107,118
rows, including 10,626 distinct transfer-event tuples.  Second, the synthetic
motif generator reused attacker addresses across instances; correcting it
changes 82 historical attack senders to 689 and invalidates every reported
motif-performance metric.  A separate synthetic benchmark lacks its generator,
row-level scores, and benchmark log and is withdrawn.  None of these artifacts
supports fraud accuracy, banking generalization, adversarial robustness, or an
operational threshold.

We replace the easy-to-vary claim that agents are intrinsically fraudulent
with a causal explanation: loss becomes possible when the chain from an
authenticated controller and bounded mandate, through instruction provenance
and tool/wallet authorization, to settlement is not preserved.  Automation
changes speed, scale, delegation depth, and coordination; it does not turn
actor type into fraud.  This revision contributes an explicit evidence vector,
threat-role and observation-plane model, claim ledger, fail-closed research
score, and preregisterable validation design.  Stronger conclusions require an
immutable event-complete dataset, verified execution and control evidence,
independently adjudicated financial-crime outcomes, leakage-safe temporal or
component holdouts, and external replication.

\medskip
\noindent\textbf{Keywords:} agentic payments, construct validity, fraud
detection, ERC-8004, transaction graphs, research integrity
\end{abstract}
